\section{Three-Layer Semantic Storage Architecture}

\label{sec:architecture}

DSO employs a three-layer architecture optimizing for different properties at
each level.

\subsection{Layer 1: On-Chain Registry (Integrity)}

The ExperienceRegistry smart contract stores metadata and integrity proofs for
each submitted prior. No prior data is stored on-chain; only cryptographic
commitments.

\begin{definition}[Registry Entry]
A registry entry $\mathcal{R}_i$ consists of:
\begin{equation}
    \mathcal{R}_i = \left(\texttt{merkleRoot}_i, \; \texttt{cid}_i, \;
    \texttt{schema}_i, \; q_i, \; \texttt{addr}_i, \; D_i, \; t_i\right),
\end{equation}
where $\texttt{merkleRoot}_i$ is the Merkle root of the prior data,
$\texttt{cid}_i$ is the content identifier on IPFS/Arweave, $\texttt{schema}_i$
identifies the decomposition schema, $q_i$ is the quality score, $\texttt{addr}_i$
is the submitter's address, $D_i$ is the bond amount, and $t_i$ is the
submission timestamp.
\end{definition}

\paragraph{Contract Interface.}
\begin{verbatim}
interface IExperienceRegistry {
    struct Entry {
        bytes32 merkleRoot;
        string  cid;             // IPFS or Arweave content ID
        uint64  schema;          // Schema version identifier
        uint64  quality;         // Quantized quality [0, 1000]
        address submitter;       // Submitter address
        uint256 bond;            // Staked bond amount
        uint64  timestamp;       // Submission block timestamp
    }

    function submit(Entry calldata e)
        external payable returns (uint256 id);

    function voteQuality(uint256 id, uint64 score)
        external;

    function challenge(uint256 id, bytes calldata proof)
        external returns (bool slashed);

    function aggregate(uint256[] calldata ids)
        external returns (bytes32 aggregateMerkleRoot);

    function getEntry(uint256 id)
        external view returns (Entry memory);

    function getEntriesBySchema(uint64 schema)
        external view returns (uint256[] memory ids);
}
\end{verbatim}

\subsection{Layer 2: Content-Addressed Storage (Availability)}

Prior data is stored on IPFS~\cite{benet2014ipfs} or
Arweave~\cite{williams2019arweave} for persistent availability. Each prior is
serialized as a Protocol Buffer message:

\begin{verbatim}
message ExperiencePrior {
    bytes   delta_signs = 1;     // Packed 1-bit signs
    float   scale = 2;          // Per-expert scale factor
    float   quality = 3;        // Quality score
    bytes   context_predicate = 4; // Serialized matching fn
    bytes   task_embedding = 5;  // d-dimensional embedding
    uint64  vocab_size = 6;     // |V_m|
    repeated uint32 vocab_indices = 7; // Sparse support
    bytes   merkle_proof = 8;   // Proof of inclusion
}
\end{verbatim}

The content identifier (CID) is the SHA-256 hash of the serialized message,
ensuring deduplication and integrity verification.

\subsection{Layer 3: Local Replica Database (Performance)}

Each node maintains a local LanceDB~\cite{lancedb2024} instance containing
replicas of high-quality priors. The local database supports:

\begin{itemize}[leftmargin=1.5em]
    \item \textbf{Vector search:} Nearest-neighbor queries over task embeddings
          for prior retrieval in $O(\log N)$ time.
    \item \textbf{CRDT semantics:} Conflict-free merge of priors from multiple
          sources using last-writer-wins with quality-based tie-breaking.
    \item \textbf{Garbage collection:} Automatic eviction of priors with quality
          below $q_{\min}$ or age exceeding $T_{\max}$.
    \item \textbf{Offline operation:} Full functionality without network access
          using cached priors.
\end{itemize}

\begin{definition}[Replica Consistency]
A local replica $L_i$ is \emph{$\tau$-consistent} with the registry if for
every entry $\mathcal{R}_j$ with $q_j \geq q_{\text{fetch}}$ and
$t_j \leq t_{\text{now}} - \tau$, the corresponding prior data is present in
$L_i$ and its Merkle proof verifies against $\texttt{merkleRoot}_j$.
\end{definition}

\subsection{Cross-Layer Integrity}

The three layers are connected by cryptographic commitments:

\begin{proposition}[End-to-End Integrity]
For any prior with registry entry $\mathcal{R}_i$, a verifier can confirm that
the local replica data $E_i$ is authentic by:
\begin{enumerate}[leftmargin=1.5em]
    \item Computing the Merkle root of $E_i$ and comparing with
          $\texttt{merkleRoot}_i$ from the on-chain registry.
    \item Verifying the CID of the serialized data matches $\texttt{cid}_i$.
    \item Checking the Merkle proof within the serialized data.
\end{enumerate}
Any modification to the prior data is detected with probability
$1 - 2^{-256}$ (collision resistance of SHA-256).
\end{proposition}

% ============================================================================
% 4. DSO PROTOCOL SPECIFICATION
% ============================================================================
